% CAPÍTULO 3 — MATERIAIS E MÉTODOS

\chapter{Materiais e Métodos}
\label{cap:metodologia}

Este capítulo descreve o método de pesquisa, as fontes de dados, o recorte adotado, a arquitetura efetivamente implementada, as ferramentas utilizadas e os modelos aplicados. A \autoref{sec:revisoes} declara, de forma explícita, as revisões feitas em relação à proposta apresentada no TCC~1.

\section{Método}
\label{sec:metodo}

O trabalho seguiu o CRISP-DM (\textit{Cross-Industry Standard Process for Data Mining}), proposto por \citeonline{Chapman2000} como processo padrão para projetos de mineração e análise de dados. As seis fases do método são organizadas em um ciclo iterativo, em que é possível voltar para etapas anteriores sempre que novas descobertas justifiquem revisão. A \autoref{fig:crispdm-ciclo} ilustra esse fluxo.

A natureza iterativa do método não foi um detalhe formal: foi o modo como o trabalho de fato andou. A proposta definiu as fontes, o recorte e as perguntas iniciais; a execução mostrou, repetidas vezes, que algo precisava mudar. A cada descoberta, o trabalho voltou a uma fase anterior do ciclo --- para buscar dados que não estavam previstos, corrigir a preparação, trocar um modelo ou reformular uma pergunta --- e só então seguiu adiante. O \autoref{quadro:crispdm-retornos}, ao fim desta seção, relaciona os principais retornos, com a descoberta que motivou cada um e a fase a que o trabalho voltou.

\begin{figure}[htbp]
    \centering
    \caption{Fases do ciclo de vida CRISP-DM}
    \label{fig:crispdm-ciclo}
    \begin{tikzpicture}[scale=0.8, transform shape,
        every node/.style={font=\footnotesize},
        fase/.style={ellipse, draw=black!70, thick, fill=blue!8,
                     minimum width=2.8cm, minimum height=1.0cm, align=center},
        centro/.style={circle, draw=black!70, thick, fill=orange!20,
                       minimum size=1.6cm, align=center, font=\small\bfseries},
        seta/.style={-{Stealth[length=2mm]}, thick, black!70},
        bidir/.style={{Stealth[length=2mm]}-{Stealth[length=2mm]}, thick, black!70}
    ]

    % Posicoes em circulo (raio ~4cm)
    \node[fase] (negocio)  at ( 90:4cm) {Entendimento \\ do Negócio};
    \node[fase] (dados)    at ( 30:4cm) {Entendimento \\ dos Dados};
    \node[fase] (prep)     at (-30:4cm) {Preparação \\ dos Dados};
    \node[fase] (modelo)   at (-90:4cm) {Modelagem};
    \node[fase] (avaliacao)at (210:4cm) {Avaliação};
    \node[fase] (implant)  at (150:4cm) {Implantação};

    % Centro
    \node[centro] (dataset) at (0,0) {Dados};

    % Setas entre fases (sentido horario)
    \draw[bidir] (negocio)  -- (dados);
    \draw[seta]  (dados)    -- (prep);
    \draw[bidir] (prep)     -- (modelo);
    \draw[seta]  (modelo)   -- (avaliacao);
    \draw[seta]  (avaliacao)-- (implant);
    \draw[seta]  (avaliacao) to[bend right=25] (negocio);

    % Seta circular externa indicando iteracao (raio maior que o dos balões)
    \draw[seta, dashed, black!50]
        (110:5.8cm) arc[start angle=110, end angle=-245, radius=5.8cm];

    \end{tikzpicture}
    \\[4pt]
    Fonte: Autoria própria, adaptado de \citeonline{Chapman2000}.
\end{figure}

A \autoref{tab:crispdm-fases} sintetiza o mapeamento entre as fases do CRISP-DM e as etapas efetivamente executadas.

\begin{table}[htbp]
    \centering
    \caption{Mapeamento entre as fases do CRISP-DM e as etapas executadas}
    \label{tab:crispdm-fases}
    \renewcommand{\arraystretch}{1.35}
    \begin{tabular}{p{4.2cm}p{9.8cm}}
        \toprule
        \rowcolor{black!8}
        \textbf{Fase CRISP-DM} & \textbf{Execução no trabalho} \\
        \midrule
        \textbf{Entendimento do Negócio}  & Caracterização do mercado de TI brasileiro e formulação das cinco perguntas de pesquisa da \autoref{quadro:perguntas}, que passaram a organizar tanto os resultados quanto a interface. \\
        \textbf{Entendimento dos Dados}   & Levantamento dos \textit{layouts} de cada ano, extração dos dicionários oficiais a partir das planilhas de \textit{layout} do servidor público e auditoria da tradução dos códigos; identificação dos defeitos de fonte relatados na \autoref{sec:achados-qualidade}. \\
        \textbf{Preparação dos Dados}     & Construção das camadas Bronze, Silver (restrita ao recorte de TI) e Gold; definição e posterior revisão do recorte de tecnologia; harmonização de nomes de coluna, tipagem de campos numéricos e agrupamento dos arquivos por partição. \\
        \textbf{Modelagem}                & Modelos SARIMA sobre o saldo mensal; estimador de \textit{nowcast} do estoque; tábua de vida e modelo de risco em tempo discreto para a permanência no vínculo; decomposição de Oaxaca-Blinder do hiato salarial; agrupamento de municípios por $k$-médias. \\
        \textbf{Avaliação}                & Validação de origem móvel com MAE e RMSE contra \textit{baseline} ingênuo sazonal; validação do \textit{nowcast} com um ano retido do treinamento; verificação de reprodutibilidade confrontando o dado bruto com a camada publicada. \\
        \textbf{Implantação}              & Migração do \textit{data lake} e do processamento para a nuvem, com execução automatizada e incremental; implantação do \textit{dashboard} consumindo o próprio \textit{data lake} público; publicação do código-fonte em repositório público. \\
        \bottomrule
    \end{tabular}
    \\[4pt]
    Fonte: Autoria própria, adaptado de \citeonline{Chapman2000}.
\end{table}

O mapeamento acima lê-se como uma sequência, mas o trabalho não a percorreu uma vez só. Das doze idas e vindas do \autoref{quadro:crispdm-retornos}, oito partiram da Avaliação: foi ao confrontar um resultado com o que era plausível, ou com uma segunda conta do mesmo número, que os defeitos apareceram. Quatro delas levaram de volta ao Entendimento do Negócio, isto é, mudaram o que se perguntava --- e três trouxeram fontes de dados que a proposta não previa.

\begin{quadro}[htbp]
    \centering
    \caption{Retornos no ciclo CRISP-DM motivados por descobertas durante o trabalho}
    \label{quadro:crispdm-retornos}
    \footnotesize
    \renewcommand{\arraystretch}{1.25}
    \begin{tabularx}{\textwidth}{%
        >{\raggedright\arraybackslash\hsize=1.20\hsize}X
        >{\raggedright\arraybackslash\hsize=0.62\hsize}X
        >{\raggedright\arraybackslash\hsize=0.62\hsize}X
        >{\raggedright\arraybackslash\hsize=1.56\hsize}X}
        \toprule
        \rowcolor{black!8}
        \textbf{Descoberta} & \textbf{Fase em que surgiu} & \textbf{Retorno a} & \textbf{O que mudou} \\
        \midrule
        Os microdados chegam só com códigos; as descrições estão em planilhas de \textit{layout} & Entendimento dos Dados & Preparação dos Dados & Nova fonte: dicionários extraídos das planilhas oficiais, com procedência (\autoref{subsec:dicionarios}) \\
        A série do CAGED não fecha só com as movimentações do mês & Entendimento dos Dados & Entendimento dos Dados & Novas fontes: declarações fora do prazo, exclusões e ajustes, não previstas na proposta \\
        Códigos com zero à esquerda no microdado e sem ele na planilha: tradução vazia, sem erro & Avaliação & Preparação dos Dados & Comparação numérica das chaves na tradução \\
        Ocupações de dados fora do recorte e cinco códigos sem aderência ao escopo & Avaliação & Entendimento do Negócio & Recorte revisto e camada tratada reprocessada (\autoref{subsec:revisao-recorte}) \\
        Estoque de 2022 abaixo do de 2021 e do de 2023 & Avaliação & Preparação dos Dados & Pedaços reconstruídos; verificação de completude por pedaço e recontagem a partir do dado bruto \\
        Permanência mediana de 48 anos pelo Kaplan-Meier & Avaliação & Modelagem & Tábua de vida de período e modelo de risco em tempo discreto \\
        A mesma mediana dava 3,55 ou 3,75 conforme a tabela de origem & Avaliação & Preparação dos Dados & Mediana exata e distribuição da remuneração em faixas na Gold \\
        O intervalo de predição cruza o zero em cerca de um ano & Avaliação & Entendimento do Negócio & Pergunta 5 reformulada: horizonte de 18 meses e \textit{nowcast} do estoque corrente \\
        São Paulo concentra 41\% do estoque & Avaliação & Entendimento do Negócio & Nova pergunta: as conclusões dependem de um estado? Análise de sensibilidade e filtro de território \\
        Faltava confronto externo dos totais & Avaliação & Entendimento dos Dados & Nova fonte: painel público do MTE, consultado com os mesmos códigos (\autoref{sec:validacao-mte}) \\
        \textit{Dashboard} organizado por base mostrava gráficos sem conclusão & Implantação & Entendimento do Negócio & Interface reorganizada pelas cinco perguntas de pesquisa \\
        Disco local esgotado e dependência de uma máquina ligada & Implantação & Preparação dos Dados & Migração do \textit{data lake} e do processamento para a nuvem (\autoref{subsec:migracao}) \\
        \bottomrule
    \end{tabularx}
    \\[4pt]
    Fonte: Autoria própria.
\end{quadro}

\section{Fontes de Dados}

A base do trabalho é formada pelos microdados da RAIS e do CAGED, divulgados pelo MTE, obtidos diretamente do servidor FTP público do Programa de Disseminação das Estatísticas do Trabalho (\nolinkurl{ftp://ftp.mtps.gov.br/pdet/microdados/}) e gravados na camada Bronze sem nenhuma alteração de conteúdo. As análises cobrem a RAIS de 2007 a 2025 e o CAGED de 2007 a agosto de 2026. Nesse período, o servidor público distribui 880 arquivos compactados, que somam 53,4~GB; descompactados, são 661~GB de texto, em largura fixa ou CSV. A camada Bronze guarda esse mesmo conteúdo em Apache Parquet com compressão ZSTD de nível 3, em 56,4~GB --- 11,7 vezes menos que o texto e quase o tamanho dos pacotes compactados de origem, com a diferença de que o Parquet é lido coluna a coluna, sem descompactar o arquivo inteiro. O tamanho descompactado foi medido sem baixar os arquivos: o formato 7z registra o tamanho original de cada arquivo no cabeçalho, ao fim do pacote, e basta ler esse trecho. A \autoref{fig:fluxo-dado}, na \autoref{sec:arquitetura}, mostra o volume em cada etapa, do servidor público ao navegador.


\subsection{Desafios de Engenharia da RAIS}

A RAIS muda de esquema com frequência. No intervalo processado, o número de colunas varia de 47 a 67, refletindo portarias que introduziram novas variáveis ao longo do tempo: indicadores de trabalho intermitente e parcial (a partir de 2017), tipo de salário contratual (em 2018) e o sufixo \textit{``-- código''} nas colunas codificadas (a partir de 2023). No nível do arquivo, as bases chegam compactadas em \texttt{.7z}, com um \texttt{.txt} de largura fixa em seu interior; as edições de 2023 a 2025 trazem o arquivo de dados com a extensão \texttt{.COMT}, renomeada antes da ingestão.

Um ponto não previsto na proposta apareceu aqui: a partir de 2023 os nomes de coluna deixam de ser estáveis mesmo dentro de um mesmo ano, e um arquivo de uma região pode nomear o mesmo campo de forma diferente de outro. Foi necessário construir uma rotina de harmonização de nomes, com tabela de equivalências explícita, aplicada sobre a camada já construída.

\subsection{Desafios de Engenharia do CAGED}

Entre os dois formatos do CAGED, descritos na \autoref{subsec:caged}, mudam também o vocabulário dos códigos e a própria data de referência da movimentação, o que impede tratar a série como uma só sem harmonização explícita.

Há ainda três tabelas auxiliares do CAGED que a proposta não previa e que se mostraram necessárias para fechar a série: as movimentações fora do prazo (\textit{caged\_for}), as excluídas (\textit{caged\_exc}) e os ajustes (\textit{caged\_ajustes}). Esta última tem uma particularidade relevante: a competência da movimentação vem declarada na linha, e não no nome do arquivo, de modo que um arquivo de 2010 contém correções referentes a meses de anos anteriores.

\subsection{Dicionários Oficiais e a Tradução dos Códigos}
\label{subsec:dicionarios}

Os microdados são quase inteiramente compostos de códigos: sexo, grau de instrução, raça/cor, ocupação, atividade econômica e município chegam como números sem descrição. O MTE publica essas correspondências em planilhas de \textit{layout}, no mesmo servidor dos dados.

Em vez de transcrever essas correspondências para o código-fonte, o trabalho extraiu os dicionários das planilhas oficiais e os materializou como uma tabela do próprio \textit{data lake}, com uma linha por código e as colunas de procedência: de qual planilha, de qual aba, de qual caminho no servidor e em que data a extração foi feita. A tradução na camada Silver é um \textit{LEFT JOIN} contra essa tabela, e mantém o código original ao lado da descrição --- nada é substituído.

Duas decisões dessa etapa merecem registro, porque a alternativa falha em silêncio. A primeira é que o dicionário é indexado \textbf{por tabela}: o mesmo código significa coisas diferentes em gerações diferentes do CAGED, e um \textit{join} sem essa restrição duplica linhas e infla as contagens. A segunda é que a comparação de chaves é numérica, e não textual: o microdado da RAIS traz códigos com zero à esquerda (\texttt{'01'}) e a planilha de \textit{layout} traz o mesmo código sem ele (\texttt{'1'}). Comparados como texto, apenas os códigos de dois dígitos casam, e a tradução dos demais sai vazia sem gerar erro algum.

Na arquitetura em nuvem, as planilhas convertidas ficam em um repositório próprio, separado dos microdados. O \textit{pipeline} verifica todo mês se o servidor público passou a ter planilhas de \textit{layout} ainda não extraídas e avisa quando isso ocorre, mas não as incorpora sozinho: a extração é executada manualmente, com conferência do resultado. A assimetria é deliberada. Um mês de microdados que deixe de ser processado é uma ausência visível; um dicionário substituído sem conferência produz traduções erradas sem erro algum, e esse é o modo de falha mais silencioso do projeto.

\section{O Recorte de Tecnologia}
\label{sec:recorte}

Nem a RAIS nem o CAGED possuem um campo que identifique um vínculo como sendo de tecnologia. O recorte é, portanto, uma decisão metodológica, e foi construído como a \textbf{união} de dois critérios:

\begin{itemize}
    \item \textbf{lente de setor}: o estabelecimento tem CNAE de atividade de TI (desenvolvimento de \textit{software}, serviços de TI, tratamento de dados e portais);
    \item \textbf{lente de ocupação}: o vínculo tem CBO de TI, selecionada por \textbf{família} ocupacional (os quatro primeiros dígitos), mais três códigos avulsos, e menos os códigos que a família traz sem serem de tecnologia.
\end{itemize}

A opção pela união, e não pela interseção, é o que permite responder à segunda pergunta de pesquisa: as duas lentes são preservadas como rótulos em cada linha da camada tratada, de modo que é possível separar quem exerce ocupação de TI fora de empresas de TI de quem exerce outra ocupação dentro delas. Fosse interseção, essa distinção desapareceria do dado.

A seleção por família, e não por código individual, tem uma razão prática: o MTE cria códigos novos dentro de famílias existentes ao longo do tempo, e uma lista fechada de códigos envelheceria em silêncio. O custo dessa escolha é trazer, junto, ocupações que compartilham a família sem serem de tecnologia --- e esse custo foi medido e corrigido.

\subsection{Revisão do Recorte}
\label{subsec:revisao-recorte}

O TCC~1 previa que a relação definitiva de códigos seria produzida na fase de Entendimento dos Dados e consolidada no TCC~2. Foi o que ocorreu, em duas direções.

Entraram os quatro códigos da família \textbf{2112 (estatísticos)}. A inclusão foi discutida com a orientação e se sustenta no dado: no CAGED de 2025, 72\% das admissões nesses códigos ocorreram fora de empresas de TI, e são códigos sob os quais analistas e cientistas de dados são efetivamente registrados no mercado brasileiro. Deixá-los fora significaria excluir do recorte justamente a ocupação que mais se aproxima do trabalho com dados.

Saíram cinco códigos que a seleção por família trazia sem aderência ao escopo: os pesquisadores em física, química, matemática e ciências da terra (família 2031) e o programador de máquinas-ferramenta com comando numérico (3171-15), que é da indústria. A exclusão incide apenas sobre a lente de ocupação: um pesquisador em física empregado em uma empresa de TI continua no recorte pela lente de setor, que é a informação correta.

A família 2031 permaneceu representada por um único código, \textbf{203105 (pesquisador em ciências da computação e informática)}, que é a razão de ela estar na lista desde o início. A relação completa dos códigos, com a descrição oficial e o estoque de cada um, é apresentada no \autoref{apend:cbos} e foi gerada a partir do próprio dado, e não digitada: a tabela do apêndice é um agregado da camada Gold, de modo que qualquer revisão futura do recorte se reflete nela automaticamente.

\section{Arquitetura Implementada}
\label{sec:arquitetura}

A plataforma foi construída primeiro sobre infraestrutura local --- um \textit{data lake} em MinIO na máquina de desenvolvimento, com as etapas disparadas manualmente --- e migrada, na fase final do trabalho, para uma arquitetura executada inteiramente em nuvem. Esta seção descreve a arquitetura final; os motivos da migração estão na \autoref{subsec:migracao}. A \autoref{fig:fluxo-dado} mostra os componentes e o caminho do dado entre eles.

\begin{figure}[htbp]
    \centering
    \caption{Fluxo do dado, do servidor público ao navegador, com o volume em cada etapa}
    \label{fig:fluxo-dado}
    \resizebox{\textwidth}{!}{%
    \begin{tikzpicture}[
        etapa/.style={rectangle, rounded corners=2pt, draw=black!70, thick,
                      text width=3.05cm, minimum height=19mm, align=center, font=\small,
                      execute at begin node={\hyphenpenalty=10000\exhyphenpenalty=10000}},
        publicado/.style={double, double distance=1.2pt},
        seta/.style={-{Stealth[length=2.2mm]}, thick, black!70},
        tracejada/.style={-{Stealth[length=2.2mm]}, thick, black!55, dashed},
        acao/.style={font=\footnotesize\itshape, fill=white, inner sep=1.5pt, align=center}
    ]
    % linha 1: da fonte à camada tratada
    \node[etapa, fill=blue!8] (ftp) at (0, 0)
        {Servidor FTP do MTE\\{\footnotesize\ttfamily .7z / .zip}\\[2pt]{\bfseries53,4~GB compactados}};
    \node[etapa, fill=blue!3] (txt) at (4.3, 0)
        {Microdado em texto\\{\footnotesize largura fixa ou CSV}\\[2pt]{\bfseries$\approx$ 661~GB}};
    \node[etapa, fill=orange!20, publicado] (bronze) at (8.6, 0)
        {Bronze\\{\footnotesize Parquet, ZSTD nível 3}\\[2pt]{\bfseries56,4~GB}};
    \node[etapa, fill=gray!20, publicado] (silver) at (12.9, 0)
        {Silver de TI\\{\footnotesize recorte CNAE $\cup$ CBO, traduzido}\\[2pt]{\bfseries1,3~GB}};

    % linha 2: dicionários, agregação e apresentação
    \node[etapa, fill=blue!3] (layout) at (0, -3.4)
        {Planilhas de \textit{layout}\\{\footnotesize do MTE (\ttfamily .xls)}};
    \node[etapa, fill=blue!10, publicado] (dic) at (4.3, -3.4)
        {Dicionários\\{\footnotesize\ttfamily mte-dicionarios}\\{\footnotesize código e procedência}};
    \node[etapa, fill=red!6] (dash) at (8.6, -3.4)
        {\textit{Dashboard}\\{\footnotesize React + DuckDB-Wasm, Cloudflare Pages}};
    \node[etapa, fill=yellow!25, publicado] (gold) at (12.9, -3.4)
        {Gold\\{\footnotesize 29 agregados e modelos}\\[2pt]{\bfseries< 10~MB}};

    \node[etapa, fill=white, draw=black!50, dashed, minimum height=12mm] (mod) at (12.9, -6.2)
        {{\footnotesize SARIMA, \textit{nowcast}, Oaxaca-Blinder, $k$-médias}};

    \draw[seta] (ftp) -- node[acao, above=1mm] {7-Zip} (txt);
    \draw[seta] (txt) -- node[acao, above=1mm] {DuckDB} (bronze);
    \draw[seta] (bronze) -- node[acao, above=1mm] {SQL} (silver);
    \draw[seta] (silver) -- node[acao, right=1mm] {SQL, Python} (gold);
    \draw[seta] (gold) -- node[acao, above=1mm] {\textit{build}} (dash);
    \draw[tracejada] (layout) -- node[acao, above=1mm] {extração} (dic);
    \draw[tracejada] (dic.north east) -- node[acao, sloped, above=0.5mm] {\textit{LEFT JOIN}} (silver.south west);
    \draw[tracejada] (mod) -- (gold);
    \end{tikzpicture}%
    }
    \\[4pt]
    {\footnotesize Contorno duplo: conjunto de dados público no Hugging Face Hub. Todas as etapas, do \textit{download} à Gold, rodam em executores do GitHub Actions; o \textit{dashboard} é publicado na Cloudflare Pages.}
    \\[2pt]
    Fonte: Autoria própria.
\end{figure}

Duas diferenças estruturais separam esta arquitetura da implementação local que a precedeu.

A primeira é que armazenamento e publicação deixaram de ser etapas distintas. Na versão local, o \textit{data lake} ficava em MinIO e uma rotina copiava as camadas tratadas para um repositório público de dados; na versão final, o repositório público \textbf{é} o \textit{data lake}. Cada camada da arquitetura Medalhão é um conjunto de dados no Hugging Face Hub --- seis ao todo, relacionados no \autoref{quadro:repositorios} ---, e cada etapa lê a camada anterior e grava a seguinte diretamente nele.

\begin{quadro}[htbp]
    \centering
    \caption{Conjuntos de dados que compõem o \textit{data lake}}
    \label{quadro:repositorios}
    \footnotesize
    \renewcommand{\arraystretch}{1.3}
    \begin{tabularx}{\textwidth}{%
        >{\raggedright\arraybackslash\hsize=0.95\hsize}X
        >{\raggedright\arraybackslash\hsize=0.45\hsize}X
        >{\raggedright\arraybackslash\hsize=1.60\hsize}X}
        \toprule
        \rowcolor{black!8}
        \textbf{Repositório} & \textbf{Camada} & \textbf{Conteúdo} \\
        \midrule
        \texttt{bronze-caged} & Bronze & Microdados das cinco tabelas do CAGED, convertidos para Parquet sem alteração de conteúdo, e a tabela de procedência dos dicionários. \\
        \texttt{bronze-rais} & Bronze & Microdados de vínculos e de estabelecimentos da RAIS, nas mesmas condições. \\
        \texttt{mte-dicionarios} & Referência & Planilhas de \textit{layout} do MTE convertidas para Parquet, uma por aba, usadas na tradução dos códigos. \\
        \texttt{silver-caged-ti} & Silver & CAGED traduzido e restrito ao recorte de tecnologia, particionado por ano e mês; e, na raiz, a versão consolidada em um arquivo por tabela e ano, para consulta por HTTP. \\
        \texttt{silver-rais-ti} & Silver & RAIS traduzida e restrita ao recorte de tecnologia, particionada por ano. \\
        \texttt{gold-analise-mercado-ti} & Gold & Agregados que o \textit{dashboard} consome e resultados materializados dos modelos. \\
        \bottomrule
    \end{tabularx}
    \\[4pt]
    Fonte: Autoria própria.
\end{quadro}

A segunda é que o processamento não tem mais máquina fixa. Cada etapa roda em um executor efêmero do GitHub Actions, que é criado para a execução, instala as dependências, processa e é descartado. Nada persiste entre execuções além do que foi gravado no repositório de dados --- o que, como se verá na \autoref{subsec:orquestracao-pipeline}, determinou a forma de dividir o trabalho.

A cadeia de camadas é Bronze, Silver de TI e Gold. A Silver não traduz o mercado inteiro para depois recortá-lo: como o recorte é definido de antemão pelos códigos de CBO e de CNAE, a tradução é aplicada só às linhas que pertencem a ele, diretamente a partir da Bronze. A alternativa --- uma Silver de mercado completo, da qual se extrairia a área de interesse --- teria volume da mesma ordem da Bronze, ou maior, já que a tradução acrescenta uma coluna de descrição ao lado de cada código, para servir a uma análise que consome 1,3~GB. Analisar outra área não exige essa camada intermediária: basta construir outra Silver a partir da mesma Bronze, com outro recorte.

\subsection{Orquestração do \textit{Pipeline}}
\label{subsec:orquestracao-pipeline}

O \textit{pipeline} é dividido em nove etapas, cada uma um \textit{workflow} próprio do GitHub Actions, com uma única responsabilidade. Um décimo \textit{workflow}, o orquestrador, encadeia as nove pela dependência real entre elas, e não por horários escalonados: uma etapa só começa quando as etapas de que depende terminaram. A \autoref{fig:actions}, adiante, mostra esse grafo em uma execução real.

O orquestrador roda toda semana e também pode ser disparado manualmente a partir de qualquer camada --- só a Silver e a Gold, quando a Bronze já está em dia, ou só a Gold, quando se quer apenas reagregar. Cada etapa, por sua vez, continua podendo ser executada isoladamente, o que permite testar ou repetir uma peça sem disparar as demais.

\paragraph{Processamento incremental.} As etapas de Bronze e de Silver têm duas fases. A primeira, de verificação, compara o que existe na origem com o que já está publicado e produz a lista de pares (tabela, ano) pendentes; na Silver, essa comparação é feita só pela listagem dos arquivos, sem ler nenhum dado. A segunda fase processa cada par pendente em um \textit{job} próprio, até quatro em paralelo. Por padrão, a verificação olha o ano corrente e o anterior, que concentram a publicação de dado novo pelo MTE; anos mais antigos são reprocessados sob demanda, informando o intervalo na execução manual. O particionamento por ano, descrito adiante, é o que torna essa unidade de trabalho possível: reprocessar um ano não toca nos demais.

\paragraph{Identidade da origem.} Para decidir se um arquivo do servidor público já foi ingerido, a verificação da Bronze não compara nomes de arquivo, e sim o caminho de origem gravado em cada linha da camada. A razão é concreta: arquivos grandes da RAIS são divididos em partes no momento da publicação, e 146 arquivos da Bronze da RAIS são partes de um original que, com o nome inteiro, não existe no repositório. Uma comparação por nome concluiria que esses originais nunca foram ingeridos e os baixaria de novo a cada execução. Quando a própria verificação falha --- por erro de rede, por exemplo ---, o item é tratado como já existente: deixar de atualizar um ano é uma falha visível e recuperável, enquanto ingerir duas vezes o mesmo arquivo duplicaria o dado sem erro algum.

\paragraph{Gold dividida por dependência.} A Gold é reconstruída inteira a cada execução, porque é pequena e barata de agregar; o custo do \textit{pipeline} está na tradução, e não na agregação. Ela foi dividida em quatro etapas pela dependência efetiva de dados de cada parte, verificada no código: os agregados da RAIS dependem só da Silver da RAIS; o território (mapa e perfis de município) depende das duas Silvers, e os perfis dependem do mapa da mesma rodada, por isso os dois ficam juntos; os consolidados do CAGED dependem só da Silver do CAGED; e os modelos dependem das duas Silvers e do estoque anual produzido pelos agregados da RAIS, que é usado pelo \textit{nowcast}. Como os executores não compartilham disco, cada etapa publica a própria parte ao terminar, e a etapa de modelos baixa do repositório o estoque anual de que precisa antes de começar. A divisão foi motivada por uma falha real: em uma das primeiras execuções em nuvem da Gold, ainda em etapa única, os consolidados do CAGED falharam depois que todas as outras partes já tinham terminado, e refazê-los exigia refazer tudo.

\paragraph{Limites do serviço.} O repositório de dados recusa requisições acima de uma taxa máxima (HTTP 429). As leituras do DuckDB passaram a ser autenticadas, o que eleva esse limite, e são repetidas até seis vezes, com espera crescente e tempo máximo de resposta --- sem ele, uma conexão que para de responder sem erro manteria o \textit{job} parado até o limite do executor.

\paragraph{Acompanhamento.} As etapas de Bronze registram, no início, os recursos do executor recebido (processadores, memória e disco), já que eles não são fixos entre execuções, e as etapas enviam notificações com o resultado --- se havia dado novo, o que foi processado e o que falhou --- por um \textit{bot} de mensagens. A interface do GitHub Actions exibe o grafo a cada execução, com o estado de cada etapa. Na execução da \autoref{fig:actions}, disparada a partir da Silver, as etapas de Bronze aparecem sem execução; a verificação da Silver do CAGED encontrou três pares pendentes --- as movimentações, as declarações fora do prazo e as exclusões do mês recém-publicado de 2026 ---, traduzidos em paralelo; a da RAIS não encontrou nada, e a tradução foi pulada. A Gold e o site foram reconstruídos em seguida, e a execução inteira levou 22 minutos.

\begin{figure}[htbp]
    \centering
    \caption{Execução do orquestrador no GitHub Actions, iniciada na camada Silver}
    \label{fig:actions}
    \resizebox{\textwidth}{!}{%
    \begin{tikzpicture}[
        job/.style={rectangle, rounded corners=3pt, draw=black!30, fill=white,
                    text width=3.8cm, minimum height=7.5mm, align=left, font=\small,
                    inner xsep=2mm,
                    execute at begin node={\hyphenpenalty=10000\hspace*{4.4mm}}},
        pulado/.style={text=black!50},
        matriz/.style={rectangle, rounded corners=3pt, draw=black!30, fill=black!3, inner sep=1.3mm},
        rotmatriz/.style={font=\scriptsize\ttfamily, text=black!70, anchor=south west, inner sep=1.5pt},
        etapa/.style={font=\small\bfseries, anchor=east, align=right, text=black!75},
        liga/.style={black!40, line width=0.7pt}
    ]
    \newcommand{\ok}[1]{%
        \fill[green!50!black] ([xshift=3.4mm]#1.west) circle (1.5mm);
        \draw[white, line width=0.9pt] ([xshift=2.75mm, yshift=0.05mm]#1.west) -- ([xshift=3.25mm, yshift=-0.5mm]#1.west) -- ([xshift=4.05mm, yshift=0.55mm]#1.west);}
    \newcommand{\pul}[1]{\draw[black!45, line width=0.7pt] ([xshift=3.4mm]#1.west) circle (1.5mm);}
    \newcommand{\dur}[1]{\hfill{\scriptsize\textcolor{black!55}{#1}}}

    % colunas a cada 4,2: 0; 4,2; 8,4; 12,6
    \node[etapa] at (-2.1,  0.0)  {Decisão};
    \node[etapa] at (-2.1, -1.4)  {Bronze\\verificação};
    \node[etapa] at (-2.1, -3.0)  {Bronze\\ingestão};
    \node[etapa] at (-2.1, -4.6)  {Silver\\verificação};
    \node[etapa] at (-2.1, -6.8)  {Silver\\tradução};
    \node[etapa] at (-2.1, -9.2)  {Gold};
    \node[etapa] at (-2.1, -10.75) {Publicação};

    \node[job] (dec) at (6.75,0) {decidir\dur{2s}}; \ok{dec}

    % Bronze: não executada nesta rodada
    \node[job, pulado] (vdi) at (2.25,-1.4) {dicionario}; \pul{vdi}
    \node[job, pulado] (vbc) at (6.75,-1.4) {bronze-caged}; \pul{vbc}
    \node[job, pulado] (vbr) at (11.25,-1.4) {bronze-rais}; \pul{vbr}
    \node[job, pulado] (abc) at (4.50,-3.0) {não executado}; \pul{abc}
    \node[job, pulado] (abr) at (9.00,-3.0) {não executado}; \pul{abr}

    % Silver
    \node[job] (vsc) at (4.50,-4.6) {silver-caged-ti\dur{30s}}; \ok{vsc}
    \node[job] (vsr) at (9.00,-4.6) {silver-rais-ti\dur{33s}};  \ok{vsr}
    \node[job] (t1) at (4.50,-6.0) {caged\_mov\dur{1m39s}}; \ok{t1}
    \node[job] (t2) at (4.50,-6.8) {caged\_exc\dur{46s}};   \ok{t2}
    \node[job] (t3) at (4.50,-7.6) {caged\_for\dur{51s}};   \ok{t3}
    \node[job, pulado] (t4) at (9.00,-6.0) {nada pendente}; \pul{t4}

    % Gold e site
    \node[job] (gc) at (0.00,-9.2) {gold-caged-\\\hspace*{4.4mm}consolidado\dur{13m55s}}; \ok{gc}
    \node[job] (gr) at (4.50,-9.2) {gold-rais\dur{6m59s}};       \ok{gr}
    \node[job] (gt) at (9.00,-9.2) {gold-territorio\dur{9m53s}}; \ok{gt}
    \node[job] (gm) at (13.50,-9.2) {gold-modelos\dur{13m50s}};   \ok{gm}
    \node[job] (st) at (6.75,-10.75) {site\dur{4s}};              \ok{st}

    \begin{scope}[on background layer]
        \node[matriz, fit=(abc)] (mbc) {};
        \node[matriz, fit=(abr)] (mbr) {};
        \node[matriz, fit=(t1)(t3)] (msc) {};
        \node[matriz, fit=(t4)] (msr) {};
        % ligações por camada, como a interface agrupa
        \draw[liga] (dec.south) -- (6.75,-0.7) (2.25,-0.7) -- (11.25,-0.7);
        \foreach \b in {vdi,vbc,vbr} \draw[liga] (\b.north |- 0,-0.7) -- (\b.north);
        \foreach \a in {vdi,vbc,vbr} \draw[liga] (\a.south) -- (\a.south |- 0,-2.15);
        \draw[liga] (2.25,-2.15) -- (11.25,-2.15);
        \foreach \b in {mbc,mbr} \draw[liga] (\b.north |- 0,-2.15) -- (\b.north);
        \foreach \a in {mbc,mbr} \draw[liga] (\a.south) -- (\a.south |- 0,-3.9);
        \draw[liga] (4.50,-3.9) -- (9.00,-3.9);
        \foreach \b in {vsc,vsr} \draw[liga] (\b.north |- 0,-3.9) -- (\b.north);
        \draw[liga] (vsc.south) -- (msc.north) (vsr.south) -- (msr.north);
        \draw[liga] (msc.south) -- (4.50,-8.4) (msr.south) -- (9.00,-8.4) (0.00,-8.4) -- (13.50,-8.4);
        \foreach \b in {gc,gr,gt,gm} \draw[liga] (\b.north |- 0,-8.4) -- (\b.north);
        \foreach \a in {gc,gr,gt,gm} \draw[liga] (\a.south) -- (\a.south |- 0,-10.05);
        \draw[liga] (0.00,-10.05) -- (13.50,-10.05) (st.north |- 0,-10.05) -- (st.north);
    \end{scope}
    \node[rotmatriz] at (mbc.north west) {Matrix: atualizar};
    \node[rotmatriz] at (mbr.north west) {Matrix: atualizar};
    \node[rotmatriz] at (msc.north west) {Matrix: traduzir (2026)};
    \node[rotmatriz] at (msr.north west) {Matrix: traduzir};

    \end{tikzpicture}%
    }
    \\[4pt]
    {\footnotesize Círculo cheio: etapa executada; círculo vazio: etapa não executada. As ligações seguem o agrupamento por camada da interface do GitHub Actions; as dependências exatas da Gold estão descritas no texto.}
    \\[2pt]
    Fonte: Autoria própria, a partir da execução de 22/10/2026 registrada no GitHub Actions.
\end{figure}

\subsection{Particionamento e Organização dos Arquivos}

Cada tabela das camadas tratadas é particionada no estilo \textit{hive}, com o ano (e o mês, quando a tabela é mensal) expresso no caminho: \nolinkurl{caged_mov/ano_particao=2025/mes_particao=6/}. Essa convenção permite ao DuckDB descartar partições inteiras pela leitura do caminho, antes de abrir qualquer arquivo.

A camada nasceu com um arquivo de saída por arquivo de entrada, o que é natural quando a ingestão é feita arquivo a arquivo, mas produziu 2.111 arquivos de menos de 2~MB cada --- 35 por ano apenas na tabela de vínculos da RAIS. O custo disso não está no espaço, e sim no número de objetos: cada arquivo é uma conexão, um rodapé de metadados e um plano de execução à parte; sobre HTTP, é uma requisição por arquivo. Uma rotina de consolidação reduziu a camada a 628 arquivos, um por partição, preservando a procedência de cada linha em colunas de linhagem (\texttt{arquivo\_fonte}, \texttt{caminho\_fonte}, \texttt{data\_ingestao} e \texttt{arquivo\_bronze}) em vez de no nome do arquivo.

\section{Auditoria e Garantia de Qualidade}
\label{sec:auditoria}

A validação prevista na proposta --- confrontar os agregados com os painéis oficiais --- foi complementada por uma rotina de auditoria automatizada, executada sobre a camada construída antes de qualquer análise. Ela aplica três verificações a cada tabela:

\begin{itemize}
    \item \textbf{completude}: todo arquivo do servidor público que contém linhas dentro do recorte tem saída na camada tratada. A identidade da origem é o caminho no servidor, registrado em cada linha, e não o nome do arquivo de saída --- que muda quando a camada é reparticionada ou consolidada;
    \item \textbf{deriva de esquema}: nenhuma coluna aparece, desaparece e reaparece ao longo da série. Uma coluna presente em 2023, ausente em 2024 e presente em 2025 indica erro de mapeamento, não mudança de \textit{layout}; colunas descontinuadas em um bloco contíguo de anos finais são toleradas, porque correspondem ao que a fonte fez;
    \item \textbf{tradução vazia}: nenhuma coluna traduzida tem apenas valores de ``não informado'' ou apenas nulos quando a origem tem códigos distintos, sintoma de chave incompatível no \textit{join} com o dicionário.
\end{itemize}

As três verificações passaram sem apontamentos na camada final. Mais relevante que isso, porém, é o que elas encontraram durante a construção, relatado na \autoref{sec:achados-qualidade}.

Os defeitos apontados durante a construção foram corrigidos por rotinas idempotentes --- harmonização de nomes, tipagem de colunas numéricas, consolidação de arquivos e aplicação da revisão do recorte ---, o que permitiu revisar o recorte por poda e complemento, sem reprocessar a camada inteira. O construtor atual da Silver incorpora as duas correções mais importantes, a harmonização de nomes entre as eras de \textit{layout} e a tipagem dos campos numéricos, de modo que um ano novo já entra correto.

\section{Ferramentas Utilizadas}

O \textit{pipeline} foi escrito em Python \cite{Python2024} e é executado pelo GitHub Actions \cite{GitHubActions2026}. O \textit{data lake} fica no Hugging Face Hub \cite{HuggingFaceDatasets2026}, e a biblioteca \textit{huggingface\_hub} publica os arquivos de cada etapa em um único \textit{commit}, para que uma etapa interrompida não deixe a camada pela metade. As consultas usam o DuckDB \cite{DuckDB2024}, que lê o Apache Parquet \cite{ApacheParquet2024} direto do repositório pelo protocolo \texttt{hf://} e devolve os resultados em Apache Arrow \cite{PyArrow2024}. Os modelos usam o \textit{statsmodels} \cite{Seabold2010} --- SARIMA, regressões da decomposição salarial e modelo de risco --- e o \textit{scikit-learn} \cite{Pedregosa2011}, no $k$-médias e na silhueta. O \textit{dashboard} é escrito em React e TypeScript \cite{React2026}, consulta os dados com o DuckDB-Wasm \cite{Kohn2022}, desenha os gráficos com o Apache ECharts \cite{Li2018} e é hospedado no Cloudflare Pages \cite{CloudflarePages2026}.

\section{Modelos Aplicados}

\subsection{Previsão do Saldo Mensal}
\label{subsec:metodo-sarima}

A série de saldo mensal de empregos em TI (admissões menos desligamentos) foi modelada pela família SARIMA, descrita na \autoref{subsec:sarima}. Foram avaliadas seis especificações, todas com componente sazonal de período 12,, e um \textit{baseline} ingênuo sazonal. A grade é deliberadamente pequena, pelo risco de sobreajuste na validação discutido na \autoref{subsec:origem-movel}.

A avaliação seguiu validação de origem móvel \cite{Tashman2000}: o modelo é treinado até um instante $T$, prevê os 12 meses seguintes, a origem avança 12 meses e o procedimento se repete, totalizando cinco dobras. O \textit{baseline} é avaliado nas mesmas dobras e com o mesmo horizonte. As métricas são MAE e RMSE, apresentadas na \autoref{subsec:metricas}.

\subsection{\textit{Nowcast} do Estoque}
\label{subsec:metodo-nowcast}

A RAIS de um ano só é publicada no ano seguinte, de modo que o estoque do ano corrente é sempre desconhecido. O CAGED, por outro lado, é mensal e já publicou parte do ano. Isso permite estimar o estoque corrente pela identidade contábil
\begin{equation}
    \widehat{E}_t = E_{t-1} + F_t,
\end{equation}
em que $E_{t-1}$ é o estoque do último ano publicado da RAIS e $F_t$ é o fluxo líquido acumulado do CAGED no ano $t$, com os meses ainda não divulgados projetados pelo modelo da \autoref{subsec:metodo-sarima}.

A identidade é aproximada, e não exata: a variação do estoque entre dois 31 de dezembro não é idêntica ao saldo do CAGED, porque as duas bases têm coberturas e prazos de declaração distintos. Por isso foram comparados quatro estimadores para a relação entre $\Delta E$ e $F$: a identidade direta, uma regressão linear de $\Delta E$ sobre $F$, a razão mediana histórica entre os dois e um estimador ingênuo que supõe estoque constante.

A validação foi feita de duas formas. A primeira é retrospectiva, aplicando cada estimador a todos os pares ano-a-ano disponíveis. A segunda --- mais exigente --- retém um ano inteiro do treinamento: os estimadores são ajustados apenas com dados até 2024 e usados para estimar o estoque de 2025, cujo valor verdadeiro já é conhecido mas não foi visto pelo ajuste.

\subsection{Permanência e Sobrevivência no Vínculo}
\label{subsec:metodo-sobrevivencia}

A permanência no emprego foi analisada a partir do tempo de vínculo declarado na RAIS, com desligamentos por iniciativa não comparável (aposentadoria, falecimento, transferência e fim de contrato de experiência, entre outros) tratados como censura.

O método adotado foi a \textbf{tábua de vida de período}, descrita na \autoref{subsec:sobrevivencia}, com sete faixas de tempo de casa (0--6 meses, 6--12 meses, 1--2, 2--3, 3--5, 5--10 e 10--20 anos). Para cada faixa calcula-se a taxa de desligamento observada no ano e, a partir delas, a probabilidade acumulada de permanência.

A escolha do método decorre de um erro corrigido durante o trabalho: o estimador de Kaplan-Meier \cite{Kaplan1958}, aplicado ao corte transversal, produziu permanência mediana de 48 anos (\autoref{sec:achados-qualidade}).

Complementarmente, um modelo de risco em tempo discreto (regressão logística do desligamento no ano sobre faixa de tempo de casa, sexo, escolaridade, área de TI, porte e região) estima o efeito de cada característica mantendo as demais constantes.

\subsection{Decomposição do Hiato Salarial}
\label{subsec:metodo-oaxaca}

As diferenças de remuneração por sexo e por raça/cor foram decompostas pelo método de Oaxaca-Blinder \cite{Oaxaca1973, Blinder1973}, apresentado na \autoref{subsec:oaxaca}. A variável dependente é o logaritmo da remuneração média em salários mínimos; as covariáveis são escolaridade, área de TI, porte do estabelecimento, região, setor (TI ou não), tempo de emprego, idade, idade ao quadrado e horas contratuais.

A referência é o vetor de coeficientes da regressão \textbf{conjunta} dos dois grupos \cite{Neumark1988}.

A decomposição foi calculada ano a ano, de 2007 a 2025, o que permite observar a trajetória das duas parcelas --- a explicada por características e a não explicada --- em lugar de uma fotografia única.

\subsection{Agrupamento de Municípios}
\label{subsec:metodo-clusters}

Para caracterizar a geografia do mercado além do \textit{ranking} de maiores estoques, os municípios foram agrupados por $k$-médias \cite{MacQueen1967} sobre seis variáveis padronizadas de trajetória: logaritmo do estoque, estoque por mil habitantes, crescimento no período, intensidade de fluxo, rotatividade e remuneração mediana em salários mínimos. O número de grupos foi escolhido pelo coeficiente de silhueta \cite{Rousseeuw1987}, e não fixado \textit{a priori}.

\section{Publicação dos Dados e Migração para a Nuvem}
\label{sec:publicacao}

A proposta do TCC~1 previa publicar o código-fonte e mencionava a possibilidade de publicar os dados tratados. A possibilidade virou necessidade por uma razão de implantação: o \textit{data lake} em MinIO era local, e um \textit{dashboard} hospedado em nuvem não tinha acesso a ele. O Hugging Face Hub \cite{HuggingFaceDatasets2026} foi escolhido por aceitar Parquet particionado, servir os arquivos com requisições parciais --- o que permite ao DuckDB ler só os trechos necessários --- e versionar os conjuntos, atendendo aos princípios FAIR \cite{WilkinsonFair2016}.

\subsection{Migração para a Nuvem}
\label{subsec:migracao}

Nessa primeira forma, a publicação era uma cópia do \textit{data lake} local. Três problemas da operação local levaram a fazer do repositório público o próprio \textit{data lake}. O primeiro foi de espaço: o \textit{data lake} local chegou a 72~GB, em um disco com 27~GB livres, e 56~GB eram uma duplicata da Bronze já publicada. O segundo foi de ambiente: a política de integridade de código do Windows passou a bloquear as bibliotecas nativas de \textit{scipy}, \textit{statsmodels} e \textit{scikit-learn}, e foi nos executores Linux do GitHub Actions que os modelos voltaram a ser calculados. O terceiro foi de atualização: incorporar um mês novo dependia de alguém executar as etapas na ordem certa e republicar, e enquanto isso o \textit{dashboard} mostrava dado desatualizado sem nenhum sinal. Com o orquestrador da \autoref{subsec:orquestracao-pipeline}, a verificação passou a ser semanal e automática, e a reconstrução do site é a última etapa da mesma fila.

Antes da migração, um \textit{workflow} piloto confirmou que a latência e os limites do repositório permitiam processar a Bronze a partir de um executor; depois, cada etapa foi reescrita para ler e gravar direto no repositório, e o código ligado ao MinIO foi removido. Um cuidado operacional acompanha essa escolha: o envio ao repositório apenas \textbf{acrescenta} arquivos, e uma camada reorganizada somaria as versões antigas às novas, contando linhas em dobro sem erro algum. Por isso a publicação da Silver consolida cada partição em um arquivo e remove, no mesmo \textit{commit}, tudo o que existia daquele ano; e um ano só é publicado se todos os seus arquivos de origem foram processados.

\section{Camada de Apresentação}
\label{sec:apresentacao}

O \textit{dashboard} segue a arquitetura apresentada na \autoref{sec:visualizacao}: um site estático em React \cite{React2026}, com gráficos em Apache ECharts \cite{Li2018}, publicado no Cloudflare Pages \cite{CloudflarePages2026}, em que as consultas rodam no navegador com o DuckDB-Wasm \cite{Kohn2022}. Mudar um filtro é uma consulta local, não há servidor de aplicação para manter, e o motor é o mesmo DuckDB do \textit{pipeline}: uma consulta escrita para a Gold vale igual no site.

Os dados do site são os da Gold. Para que o CAGED pudesse ser filtrado livremente no navegador, a etapa da Gold do CAGED passou a agregar a camada tratada em dois cubos, publicados junto das demais tabelas: um no grão de ano, unidade federativa, setor da empresa, lente do recorte, área de TI, sexo, raça/cor, escolaridade e faixa etária, com 840.935 linhas, e outro no grão de ano, ocupação e lente. Os cubos guardam a \textbf{soma} e a \textbf{contagem} do salário, e não a média: a divisão é feita na consulta, para que qualquer combinação de filtros produza a média ponderada correta --- a média de médias de grupos de tamanhos diferentes seria plausível e errada.

O site embarca uma \textbf{cópia} da Gold, de cerca de 8~MB, como medida de disponibilidade: ele abre e responde mesmo com o repositório de dados fora do ar. Como os dados vão dentro do site, um mês novo só aparece depois da reconstrução, que o Cloudflare Pages faz a partir do repositório do código e que a última etapa do orquestrador dispara; se o repositório de dados estiver indisponível nesse momento, a construção falha e a versão anterior continua no ar.

A organização segue o princípio de \citeonline{Shneiderman1996}. Cada página corresponde a uma pergunta de pesquisa e abre pela própria pergunta e por poucos indicadores; os filtros vêm em seguida; e cada gráfico tem a tabela com os valores exatos a um clique. As respostas e a interpretação ficam em uma página de análise à parte, a única com texto corrido, para que as páginas de gráficos funcionem como um painel. Um filtro de território vale para o site inteiro: o visitante escolhe estados, regiões ou atalhos (como ``Brasil sem SP'') em uma lista com caixas de seleção, ou clica em um estado no mapa ou no \textit{ranking}, e todas as páginas passam a mostrar o recorte. Nem toda fonte permite esse recorte --- a decomposição salarial, a tábua de sobrevivência e os modelos de previsão foram calculados para o país ---, e esses gráficos exibem, com o filtro ligado, um selo que declara que continuam nacionais, em vez de parecerem filtrados sem estar.

Um cuidado de método acompanha o filtro: mediana não se soma (\autoref{sec:achados-qualidade}). A Gold guarda a mediana exata de cada recorte usado neste texto e a distribuição da remuneração em faixas de 0,1 salário mínimo por ano, estado e lente. Sem filtro, o site mostra a mediana exata; com filtro, soma as faixas dos estados escolhidos e interpola dentro da faixa que contém o meio da distribuição, com erro máximo de 0,05 salário mínimo.

\section{Revisões em Relação à Proposta do TCC~1}
\label{sec:revisoes}

Projetos de dados raramente terminam idênticos ao que foram propostos, e declarar as mudanças é parte do método. O \autoref{quadro:revisoes} relaciona o que foi proposto, o que foi efetivamente feito e por quê.

\begin{quadro}[htbp]
    \centering
    \caption{Revisões em relação à proposta apresentada no TCC~1}
    \label{quadro:revisoes}
    \footnotesize
    \renewcommand{\arraystretch}{1.3}
    \begin{tabularx}{\textwidth}{%
        >{\raggedright\arraybackslash\hsize=0.75\hsize}X
        >{\raggedright\arraybackslash\hsize=0.95\hsize}X
        >{\raggedright\arraybackslash\hsize=1.30\hsize}X}
        \toprule
        \rowcolor{black!8}
        \textbf{Proposto no TCC~1} & \textbf{Executado} & \textbf{Motivo} \\
        \midrule
        Modelo Prophet & Família SARIMA, com seis especificações e \textit{baseline} ingênuo sazonal & A série mensal é completa, sem lacunas, e as vantagens do Prophet (tolerância a dados faltantes e \textit{changepoints} manuais) não se aplicavam. O Prophet não foi avaliado, o que fica declarado como limitação. \\
        Projeção até 2030 & Horizonte de 18 meses & O intervalo de predição cruza o zero em cerca de um ano; uma projeção até 2030 não teria conteúdo informativo (\autoref{sec:resultado-q5}). \\
        Previsão como única modelagem & Previsão, \textit{nowcast}, sobrevivência, decomposição salarial e agrupamento & As perguntas de pesquisa formuladas na fase de Entendimento do Negócio pedem mais do que extrapolação de série: qualidade do emprego, equidade e geografia exigem métodos próprios. \\
        Publicação apenas do código & Publicação do código e de seis conjuntos de dados abertos, que formam o próprio \textit{data lake} & Primeiro, requisito de implantação: o \textit{dashboard} em nuvem não alcançava o \textit{data lake} local. Depois, consequência da migração para a nuvem. Ver \autoref{sec:publicacao}. \\
        \textit{Data lake} compatível com a API S3 (MinIO), em infraestrutura local & \textit{Data lake} em repositórios públicos do Hugging Face Hub, com processamento em executores do GitHub Actions, orquestrado por dependência e incremental por ano & Espaço em disco, bloqueio das bibliotecas de modelagem e dependência de execução manual (\autoref{subsec:migracao}). \\
        Validação contra painéis oficiais & Validação contra painéis oficiais, auditoria automatizada e teste de reprodutibilidade a partir do dado bruto & A auditoria encontrou defeitos na fonte que a comparação de agregados não revelaria, por se cancelarem na soma. \\
        Recorte por CBO & Recorte pela união de CNAE e CBO, com as duas lentes preservadas & Restringir a ocupações eliminaria do escopo quem trabalha em empresa de TI em função de apoio, e impediria a comparação entre as duas lentes --- que se tornou o achado central do trabalho. \\
        \bottomrule
    \end{tabularx}
    \\[4pt]
    Fonte: Autoria própria.
\end{quadro}

\section{Limitações e Reprodutibilidade}
\label{sec:limitacoes-metodo}

A RAIS e o CAGED cobrem exclusivamente o mercado \textbf{formal}. Trabalhadores por conta própria, profissionais autônomos sem CNPJ, microempreendedores individuais e a parcela informal do setor ficam fora do escopo, o que é especialmente relevante em uma área em que a contratação como pessoa jurídica é comum. A CBO, por sua vez, é atualizada lentamente em relação à criação de ocupações no setor de tecnologia.

A localização de cada vínculo é a do estabelecimento declarante, não a do trabalho efetivo. Como empresas de TI escolhem onde se registrar por motivos tributários, municípios com incentivo fiscal aparecem com concentração acima da real. O trabalho remoto agrava a distorção, e não há no dado público informação que permita corrigi-la --- a limitação é declarada, não contornada.

A arquitetura em nuvem troca a dependência de uma máquina pela dependência de dois serviços gratuitos de terceiros. Os limites de cada um --- a taxa de requisições do repositório de dados e a memória, o disco e a duração máxima de cada execução do GitHub Actions --- não são contratuais e podem mudar, e uma mudança de política de qualquer dos dois exigiria adaptar o \textit{pipeline}. O risco é atenuado por dois fatos: os dados estão em formato aberto, transferível para qualquer armazenamento de objetos, e as etapas são programas Python comuns, que a plataforma de integração contínua apenas agenda, encadeia e executa.

A reprodutibilidade foi tratada em quatro níveis. O código-fonte completo do \textit{pipeline} e do \textit{dashboard} está publicado em repositório público. As camadas tratadas e os agregados estão publicados como dados abertos, de modo que as análises podem ser refeitas sem executar a ingestão. A própria correspondência entre o dado bruto e o dado publicado é verificável por uma consulta, apresentada na \autoref{sec:reprodutibilidade}. E o \textit{pipeline} inteiro é reexecutável por terceiros: como os \textit{workflows} estão versionados junto do código, uma cópia do repositório, com credenciais próprias e os nomes dos repositórios de destino ajustados, executa a mesma cadeia nos mesmos executores, sem instalar nada localmente.
